\ifja
\chapter{電卓の物理学：ハードウェアレジスタによる即時検証}
\label{chap:calculator_physics}

AIやクラウドが存在する時代に、なぜ机の上の物理電卓を叩くのか。それは「計算するため」ではなく、\textbf{「AIや他人が出してきた数字の違和感を、3秒で直感的に検算・破壊するため」}である。

\section{キー操作と内部バッファの挙動}
実務電卓（CASIO/SHARP系）は、CPUの内部レジスタと直結した高効率なステートマシンである。
\begin{itemize}
    \item \textbf{C（Clear）/ CE（Clear Entry）}: 入力バッファのクリア。
    \item \textbf{M+ / M-}: ディスプレイの表示値を独立アキュムレータ（$M$）に加算／減算。
    \item \textbf{RM / CM（またはMRC）}: 独立メモリの読み出し（Recall）および初期化（Clear）。
    \item \textbf{GT（Grand Total）}: 「$=$」キー押下によって確定された値を累積蓄積するシフトレジスタ。
\end{itemize}

\section{定数乗算・除算とGTの漸化式}
実務電卓における連打操作は、等比数列およびその級数のハードウェア評価である。

\begin{align}
    y_1 &= a \div b = \frac{a}{b} \\
    y_2 &= a \div b^2 = \frac{a}{b^2} \\
    &\;\;\vdots \nonumber \\
    y_n &= a \div b^n = \frac{a}{b^n} \\
    \text{GT}_n &= \sum_{i=1}^n \frac{a}{b^i}
\end{align}
\begin{equation}
    PV_{\text{annuity}} = \sum_{i=1}^n \frac{C}{(1+r)^i}
\end{equation}
\begin{equation}
    \underbrace{\text{Assets} + \text{Expenses}}_{\ifja \text{左側（運用）: } [M+] \else \text{Left (Application): } [M+] \fi} = \underbrace{\text{Liabilities} + \text{Equity} + \text{Revenues}}_{\ifja \text{右側（調達）: } [M-] \else \text{Right (Source): } [M-] \fi}
\end{equation}


\begin{calcbox}{定数計算による現在価値（PV）と年金現価の即時算定}
将来キャッシュフロー $C$、割引率 $r$、期間 $n$ の年金現価合計：
\textbf{キー操作手順：}
\begin{enumerate}
    \item $[AC]$ でクリア後、$[C] \div [1 + r] [=]$（第1期の現在価値がGTに加算）
    \item $[=]$ を残り $(n-1)$ 回押下（各期の割引額が順次加算）
    \item $[GT]$ を押下 $\to$ 年金現価の総和が一撃で表示される。
\end{enumerate}
\end{calcbox}

\section{独立メモリによる会計方程式の解法}
左辺の要素を $[M+]$、右辺の要素を $[M-]$ に次々と放り込み、最後に $[RM]$ を押す。表示が「$0$」になれば方程式は完璧に平衡している。差額が出た場合、その値こそが「帳簿の歪み（脱漏）」である。

\else
\chapter{Physics of the Calculator: Instant Hardware Register Verification}
\label{chap:calculator_physics}

We execute physical calculations not merely to compute, but to falsify discrepancies produced by third parties or LLMs within three seconds.

\section{Key Operations and Internal Buffer Dynamics}
Commercial desktop calculators operate as direct finite-state machines over CPU registers:
\begin{itemize}
    \item \textbf{C / CE}: Clear entry and input buffer reset.
    \item \textbf{M+ / M-}: Add/subtract display value directly to independent accumulator ($M$).
    \item \textbf{RM / CM}: Recall and clear the independent memory accumulator.
    \item \textbf{GT (Grand Total)}: Accumulating shift register for committed "$=$" outputs.
\end{itemize}

\section{Constant Calculation and Recurrence Relations}
\begin{align}
    y_1 &= a \div b = \frac{a}{b} \\
    y_2 &= a \div b^2 = \frac{a}{b^2} \\
    &\;\;\vdots \nonumber \\
    y_n &= a \div b^n = \frac{a}{b^n} \\
    \text{GT}_n &= \sum_{i=1}^n \frac{a}{b^i}
\end{align}
\begin{equation}
    PV_{\text{annuity}} = \sum_{i=1}^n \frac{C}{(1+r)^i}
\end{equation}
\begin{equation}
    \underbrace{\text{Assets} + \text{Expenses}}_{\ifja \text{左側（運用）: } [M+] \else \text{Left (Application): } [M+] \fi} = \underbrace{\text{Liabilities} + \text{Equity} + \text{Revenues}}_{\ifja \text{右側（調達）: } [M-] \else \text{Right (Source): } [M-] \fi}
\end{equation}


\begin{calcbox}{Instant Present Value (PV) Evaluation via Division Recurrence}
Given recurring cash flow $C$, discount rate $r$, and horizon $n$:
\begin{enumerate}
    \item Press $[AC]$, then input $[C] \div [1 + r] [=]$ (Period 1 added to GT).
    \item Press $[=]$ for the remaining $(n-1)$ iterations.
    \item Press $[GT] \to$ The total $PV_{\text{annuity}}$ is evaluated instantly.
\end{enumerate}
\end{calcbox}

\section{Solving the Accounting Equation via Independent Accumulators}
Feed debit/application terms to $[M+]$ and credit/source terms to $[M-]$. Pressing $[RM]$ must yield exactly $0$. Any non-zero scalar explicitly reveals uncalibrated system leakage.
\fi
